% Part:sets-functions-relations
% Chapter: size-of-sets
% Section: comparing-sizes

\documentclass[../../../include/open-logic-section]{subfiles}

\begin{document}

\olfileid{sfr}{siz}{car}

\olsection{Himpunan kang Beda Gedhene, lan Teorema Cantor}

\begin{explain}
Kita wis menehi andharan kang pas babagan gagasan yen rong himpunan
padha gedhene. Kita uga bisa menehi andharan kang pas babagan gagasan
yen sawijining himpunan luwih cilik tinimbang liyane. Definisi
"luwih cilik tinimbang (utawa padha cacahe)" mbutuhake,
dudu !!a{bijection} antarane himpunan, nanging !!a{injection}
saka himpunan kapisan menyang kapindho. Yen fungsi kaya mangkono ana,
gedhene himpunan kapisan kurang saka utawa padha karo gedhene himpunan
kapindho. Kanthi intuisi, !!a{injection} saka sawijining himpunan menyang
liyane njamin yen cacah !!{element}s ing jangkauan fungsi paling ora
padha karo ing domain, amarga ora ana rong !!{element}s saka domain
kang dipetakake menyang !!{element} jangkauan kang padha.
\end{explain}

\begin{defn}
$A$ \emph{ora luwih gedhe tinimbang}~$B$, ditulis $\cardle{A}{B}$,
yen lan mung yen ana !!a{injection} $f \colon A \to B$.
\end{defn}

Cetha yen relasi iki refleksif lan transitif, nanging ora simetris
(iki dadi gladhen). Kita uga bisa ngenalake gagasan kang
ngandharake yen sawijining himpunan luwih cilik kanthi teges ketat
tinimbang liyane.

\begin{defn}
$A$ \emph{luwih cilik tinimbang}~$B$, ditulis $\cardless{A}{B}$,
yen lan mung yen ana !!a{injection}~$f\colon A \to B$ nanging ora ana
!!{bijection}~$g\colon A \to B$, yaiku $\cardle{A}{B}$ lan $\cardneq{A}{B}$.
\end{defn}

Cetha yen relasi iki irrefleksif lan transitif.
(Iki dadi gladhen.) Nganggo notasi iki, kita bisa ngandharake
yen himpunan $A$ iku !!{enumerable} yen lan mung yen $\cardle{A}{\Nat}$,
lan $A$ iku !!{nonenumerable} yen lan mung yen $\cardless{\Nat}{A}$.
Mula kita bisa ngandharake maneh
\oliflabeldef{sfr:siz:nen-alt:thm:nonenum-pownat}{%
\olref[sfr][siz][nen-alt]{thm:nonenum-pownat}
minangka pangamatan yen
$\cardless{\Nat}{\Pow{\Nat}}$}{%
\olref[sfr][siz][nen]{thm:nonenum-pownat}
minangka pangamatan yen $\cardless{\PosInt}{\Pow{\PosInt}}$}.
Sejatine, \citet{Cantor1892} mbuktekake yen asil pungkasan iki
\emph{lumaku kanthi umum}:

\begin{thm}[Cantor]\ollabel{thm:cantor}
$\cardless{A}{\Pow{A}}$, kanggo saben himpunan $A$.
\end{thm}

\begin{proof}
Pemetaan $f(x) = \{x\}$ iku !!a{injection} $f \colon A \to \Pow{A}$,
amarga yen $x \neq y$, mula $\{x\} \neq \{y\}$ miturut ekstensionalitas,
saengga $f(x) \neq f(y)$. Dadi $\cardle{A}{\Pow{A}}$.

\begin{editorial}
Kita menehi bukti kang rinci yen \olref[nen]{sec} ana ing wacan;
yen ora, diwenehake bukti kang luwih ringkes, cocog karo \olref[nen-alt]{sec}.
\end{editorial}

\oliflabeldef{sfr:siz:nen:sec}{%
  Saiki kita bakal nuduhake yen ora bisa ana fungsi !!a{surjective}~$g\colon A \to
  \Pow{A}$, luwih-luwih kang !!a{bijective}, mula
  $\cardneq{A}{\Pow{A}}$. Upamakna $g\colon A \to \Pow{A}$.
  Amarga $g$ total, saben $x \in A$ dipetakake menyang subhimpunan $g(x)
  \subseteq A$. Kita bisa nuduhake yen $g$ ora bisa surjektif.
  Carane yaiku netepake subhimpunan~$\overline{A} \subseteq A$
  kang, miturut definisine, ora bisa ana ing jangkauan~$g$. Tetepna
  \[
  \overline{A} = \Setabs{x \in A}{x \notin g(x)}.
  \]
  Amarga $g(x)$ ditetepake kanggo kabeh $x \in A$, cetha yen $\overline{A}$
  iku subhimpunan~$A$ kang ditetepake kanthi becik. Nanging himpunan iki
  ora bisa ana ing jangkauan~$g$. Upamakna $x \in A$ sembarang;
  kita bakal nuduhake yen $\overline{A} \neq g(x)$.
  Yen $x \in g(x)$, mula ora nyukupi $x \notin g(x)$, dadi
  miturut definisi~$\overline{A}$, kita duwe $x \notin
  \overline{A}$. Yen $x \in \overline{A}$, kudu nyukupi
  sipat kang netepake~$\overline{A}$, yaiku $x \in A$ lan $x \notin
  g(x)$. Amarga $x$ sembarang, iki nuduhake yen kanggo saben $x \in
  A$, $x \in g(x)$ yen lan mung yen $x \notin \overline{A}$,
  mula $g(x) \neq \overline{A}$. Domain kuantifikasi dibenerake dadi A
  amarga sumber Inggris keliru mbatesi marang himpunan diagonal (OLSIZ-007).
  Kanthi tembung liya, $\overline{A}$
  ora bisa ana ing jangkauan~$g$, kang mbantah anggepan yen~$g$
  surjektif.}{Isih kudu dituduhake yen $\cardneq{A}{\Pow{A}}$.
  Kanggo bukti kanthi kontradiksi, upamakna $\cardeq{A}{\Pow{A}}$,
  yaiku ana !!{bijection} $g \colon A \to \Pow{A}$. Saiki gatekna:
  \[
    D = \Setabs{x \in A}{x \notin g(x)}
  \]
  Gatekna yen $D \subseteq A$, mula $D \in \Pow{A}$. Amarga $g$
  iku !!a{bijection}, ana $y \in A$ saengga $g(y) = D$.
  Nanging saiki kita duwe:
  \[
    y \in g(y) \text{ yen lan mung yen } y \in D \text{ yen lan mung yen } y \notin g(y).
  \]
  Iki kontradiksi; mula $\cardneq{A}{\Pow{A}}$.}{}
\end{proof}

\begin{explain}
\oliflabeldef{sfr:siz:nen:thm:nonenum-pownat}{Mbandhingake bukti
  \olref{thm:cantor} karo bukti \olref[nen]{thm:nonenum-pownat}
  bisa menehi pangerten. Ing kana kita nuduhake yen kanggo saben dhaptar
  $Z_1$, $Z_2$, \dots, saka subhimpunan~$\PosInt$, kita bisa mbangun
  himpunan wilangan~$\overline{Z}$ kang mesthi ora ana ing dhaptar.
  Himpunan iku mesthi ora ana ing dhaptar amarga kanggo saben $n \in
  \PosInt$, $n \in Z_n$ yen lan mung yen $n \notin \overline{Z}$.
  Kanthi mangkono, mesthi ana wilangan kang dadi !!a{element} saka siji
  ing antarane $Z_n$ lan $\overline{Z}$ nanging dudu anggota liyane.
  Ing kene kita nggunakake gagasan kang padha, nanging indeks~$n$
  saiki dadi !!{element}s saka~$A$, dudu saka~$\PosInt$.
  Himpunan $\overline{A}$ ditetepake supaya beda karo~$g(x)$
  kanggo saben $x \in A$, amarga $x \in g(x)$ yen lan mung yen
  $x \notin \overline{A}$. Maneh, mesthi ana !!a{element} saka~$A$
  kang dadi !!a{element} saka siji ing antarane $g(x)$ lan $\overline{A}$
  nanging dudu anggota liyane. Kaya $\overline{Z}$ kang ora bisa ana
  ing dhaptar $Z_1$, $Z_2$, \dots, $\overline{A}$ uga ora bisa
  ana ing jangkauan~$g$.}{}

\oliflabeldef{sfr:siz:nen-alt:thm:nonenum-pownat}{Mbandhingake bukti
  \olref{thm:cantor} karo bukti \olref[nen-alt]{thm:nonenum-pownat}
  bisa menehi pangerten. Ing kana kita nuduhake yen kanggo saben dhaptar
  $N_0$, $N_1$, $N_2$, \dots, saka subhimpunan~$\Nat$, kita bisa mbangun
  himpunan wilangan~$D$ kang mesthi ora ana ing dhaptar.
  Himpunan iku mesthi ora ana ing dhaptar amarga $n \in N_n$
  yen lan mung yen $n \notin D$, kanggo saben $n \in \Nat$.
  Ing kene kita nggunakake gagasan kang padha, nanging indeks~$n$
  saiki dadi !!{element}s saka~$A$, dudu saka~$\Nat$.
  Himpunan $D$ ditetepake supaya beda karo~$g(x)$ kanggo saben $x
  \in A$, amarga $x \in g(x)$ yen lan mung yen $x \notin D$.}{}

Bukti iki uga migunani yen dibandhingake karo bukti Paradoks Russell,
\olref[sfr][set][rus]{thm:russells-paradox}. Pancen, Teorema Cantor
dadi inspirasi kanggo paradokse Russell.
\end{explain}

\begin{prob}
  Tuduhna yen ora bisa ana !!a{injection} $g\colon \Pow{A} \to
  A$, kanggo saben himpunan~$A$. Pituduh: upamakna $g\colon \Pow{A} \to A$
  iku !!{injective}. Gatekna $D = \Setabs{g(B)}{B \subseteq A \text{ lan
  } g(B) \notin B}$. Upamakna $x = g(D)$. Gunakna kasunyatan
  yen $g$ iku !!{injective} kanggo ngolehake kontradiksi.
\end{prob}

\end{document}
